% Documento principal para Overleaf. Estructura según el enunciado (sección "Documentación").
\documentclass[12pt,a4paper]{article}
\usepackage[spanish]{babel}
\usepackage[utf8]{inputenc}
\usepackage{graphicx}
\usepackage{booktabs}
\usepackage{hyperref}
\usepackage[margin=2.5cm]{geometry}

\graphicspath{{figuras/}}

\begin{document}

% 1. Portada
\begin{titlepage}
  \centering
  {\Large Instituto Tecnológico de Costa Rica\par}
  {\large Ingeniería en Computación — Campus Tecnológico Local San Carlos\par}
  \vspace{2cm}
  {\Huge\bfseries JarvisTEC\par}
  {\Large Asistente personal con reconocimiento de voz, emociones y aprendizaje automático\par}
  \vspace{2cm}
  {\large Inteligencia Artificial — I Semestre 2026\par}
  {\large Profesor: Ing. Efrén Jiménez Delgado\par}
  \vspace{1cm}
  {\large Integrantes:\par}
  % TODO: nombres y carnés
  \vfill
  {\large \today\par}
\end{titlepage}

\tableofcontents
\newpage

% 2. Solución planteada
\section{Solución planteada}
% TODO: arquitectura general (frontend, API REST, servicios de Azure/Google, modelos)

% 3. Arquitecturas de ML
\section{Arquitecturas de aprendizaje automático}
\subsection{Agente de reconocimiento de emociones}
% TODO: PEAS y tipo de agente (ver backend/features/vision_facial/SPEC.md)
\subsection{Agente de voz a texto y ejecución de comandos}
% TODO: PEAS y tipo de agente (ver backend/features/asistente_voz/SPEC.md)
\subsection{Pipeline común de los modelos de aprendizaje automático}
% TODO

% 4. Justificación y explicación de los modelos (uno por modelo, desde cada analisis.md)
\section{Justificación y explicación de los modelos}
% \input{modelos/modelo_02_autos}

% 5. Análisis de resultados
\section{Análisis de resultados obtenidos}
% TODO

% 6. Bibliografía
\bibliographystyle{ieeetr}
\bibliography{referencias}

% 7. Anexos
\appendix
\section{Anexos}
% TODO: especificación de la API, comandos de voz soportados

\end{document}
